\documentclass[]{article}

\usepackage{amsmath, amsfonts, amssymb, dsfont}
\usepackage{float}
\usepackage{fancyhdr}
\usepackage{fontawesome5}
\usepackage{mathrsfs}
\usepackage{amsmath}
\allowdisplaybreaks 
\usepackage{arydshln} 
\usepackage{tikz}
\usetikzlibrary{automata, positioning}

\usepackage[utf8]{inputenc}


\usepackage{graphicx, adjustbox, float, xcolor}
\usepackage{tikz, pgfplots}
\pgfplotsset{compat=1.17}

\usetikzlibrary{arrows.meta, automata, positioning, decorations.pathreplacing, decorations.markings}
\newcommand{\interior}[1]{%
  {\kern0pt#1}^{\mathrm{o}}%
}

\usepackage{booktabs, enumitem, caption}


\usepackage{minted} 
\usemintedstyle{friendly}
\setminted{fontsize=\small,autogobble,mathescape}

\usepackage[a4paper,
  includehead,includefoot,
  top=5mm,            
  headheight=15pt,    
  headsep=1pt,        
  bottom=20mm,left=40mm,right=40mm
]{geometry}

\pagestyle{fancy}


\lhead{Ben Vickers - 304024}
\rhead{Numerical Analysis Assignment 2}
\setlength{\headsep}{0.4cm}
\usepackage{titling}
\setlength{\droptitle}{-70pt} 

\title{\textbf{Numerical Analysis Assignment 2}}
\author{Ben Vickers - 304024}
\date{\textbf{Due: 13:59, Thursday 20 November 2025}}

\begin{document}

	\maketitle
\noindent\textbf{Assessed exercise 2.1} Consider the following functions from \(\mathbb{R}\) to \(\mathbb{R}\). Prove that \(0\):
\begin{itemize}
  \item is a simple root of the identity (also known as linear unit) \(x \mapsto x =: \mathrm{id}_{\mathbb{R}}(x)\).
  \item is not a simple root of \(x \mapsto x^{2}\).
  \item is a simple root of \(x \mapsto |x|\).
  \item is a simple root of \(x \mapsto \max\{x,2x\}\).
  \item is not a simple root of the so-called rectified linear unit (ReLU) \(x \mapsto \max\{0,x\}\).
\end{itemize}

\noindent \textbf{SOLUTION:}

\begin{enumerate}
    \item To prove that $0$ is a simple root of the identity function $f(x)=x$ we can observe that:
    \[
\displaystyle \phi(x)=\frac{\|f(x)\|}{\|x-0\|}
= \frac{\|x\|}{\|x\|}=1 > 0 \quad \text{for } x\neq 0.
\]
Hence,
\[
\displaystyle \liminf_{x\to 0}\phi(x)=\liminf_{x\to 0} 1 = 1>0
\]
so $0$ is a simple root of $f. \quad \Box$

\item To prove that $0$ is not a simple root of $f(x) = x^2$ we can observe that:
    \[
\displaystyle \phi(x)=\frac{\|f(x)\|}{\|x-0\|}
= \frac{\|x^2\|}{\|x\|}=\|x\| \quad \text{for } x\neq 0.
\]
Hence,
\[
\displaystyle \liminf_{x\to 0}\phi(x)=\liminf_{x\to 0} \|x\| =0.
\]
so $0$ is not a simple root of $f. \quad \Box$

\item To prove that $0$ is a simple root of $f(x) = |x|$ we can observe that:
    \[
\displaystyle \phi(x)=\frac{\|f(x)\|}{\|x-0\|}
= \frac{\|\,|x|\,\|}{\|x\|}=\frac{\|x\|}{\|x\|}=1 > 0 \quad \text{for } x\neq 0.
\]
Hence,
\[
\displaystyle \liminf_{x\to 0}\phi(x)=1>0
\]
so $0$ is a simple root of $f. \quad \Box$


\item To prove that $0$ is a simple root of $f(x) = \max\{x,2x\}$ we can observe that:
    \[
\displaystyle \phi(x)=\frac{\|f(x)\|}{\|x-0\|}
= \frac{\|\max\{x,2x\}\|}{\|x\|}=
\begin{cases}
\frac{\|2x\|}{\|x\|}=2 &;\, x>0\\[4pt]
\frac{\|x\|}{\|x\|}=1 & ;\,x<0
\end{cases}
\qquad > 0 \quad \text{for } x\neq 0.
\]
Hence,
\[
\displaystyle \lim_{x\to 0^{+}}\phi(x)=2 \text{ and } \lim_{x\to 0^{-}}\phi(x)=1\, \implies \displaystyle \liminf_{x\to 0}\phi(x) =1 >0.
\]
so $0$ is a simple root of $f. \quad \Box$
\item To prove that $0$ is not a simple root of $f(x)= \max\{0,x\}$ we can observe that:
    \[
\displaystyle \phi(x):=\frac{\|f(x)\|}{\|x-0\|}
= \frac{\|\max\{0,x\}\|}{\|x\|}=
\begin{cases}
\frac{\|x\|}{\|x\|}=1 &;\, x>0\\[4pt]
\frac{\|0\|}{\|x\|}=0 &;\, x<0
\end{cases} \quad \text{for } x\neq 0.
\]
Hence,
\[
\lim_{x\to 0^{+}}\phi(x)=1 \text{ and } \lim_{x\to 0^{-}}\phi(x)=0\, \implies \displaystyle \liminf_{x\to 0}\phi(x) =0.
\]
so $0$ is not a simple root of $f. \quad \Box$

\end{enumerate}

\noindent\textbf{Assessed exercise 2.2}

\begin{itemize}
  \item Check that the trapezoidal rule is not exact on polynomials of degree two.
  \item Check that the midpoint rule is exact on polynomials of degree less than or equal to 1.
\item Check that the midpoint rule is not exact on polynomials of degree two.
\end{itemize}

\noindent \textbf{SOLUTION:}

\begin{enumerate}
    \item To prove that the trapezoidal rule is not exact on polynomials of degree two, we observe that given:
    \[
    \mathcal{T}_{\,a}^{b}[f] := (b-a)\frac{f(a)+f(b)}{2}
    \]
    \noindent For polynomials of degree $2$ on the interval $[a,b]$ it holds that:
\[
f(x)\equiv x^2 \implies \displaystyle \mathcal{E}_{\mathcal{T}_{\,a}^{b}}\,[f]  = \left[\int_a^b x^2\,dx\right]-\left[\left(b-a\right)\frac{b^2+a^2}{2}\right] = \frac{b^3-a^3}{3} - \frac{(b-a)(b^2+a^2)}{2}
\]
 Exactness would require:
 \[
\frac{b^3-a^3}{3} \equiv \frac{(b-a)(b^2+a^2)}{2}
 \]
 Which clearly does not hold for all values of $a,b \in \mathbb{R}$. Hence the trapezoidal rule is not exact on polynomial functions $P$ with deg$P = 2$.\newline

\noindent Note that this result aligns with Theorem 4.7 which gives us:
\[
\displaystyle
\left[\left|\int_{a}^{b} f(x)\,dx - \mathcal{T}_{\,a}^{b}[f]\right|\right]_{f\in P_2}
\leq \frac{1}{12}\,\|f''\|_{L^\infty}\,(b-a)^{3} = \frac{\left(b-a\right)^3}{6} \neq 0 \, \, ;\, b\neq a.
\]

\item To prove that the midpoint rule is exact on polynomials of degree less than or equal to 1, we observe that given:
\[
\displaystyle \mathcal{M}_{a}^{b}[f] := (b-a)\,f\left(\frac{a+b}{2}\right)
\]
We have
\begin{itemize}
    \item For polynomials of degree 0 on the interval $[a,b]$ it holds that:
\[
f(x) \equiv 1 \implies \displaystyle \mathcal{E}_{\mathcal{M}_{a}^{b}}\,[f]   = \left[\int_a^b 1\,dx\right]-\left[\left(b-a\right)\cdot1\right] = (b-a) - (b-a) = 0.
\]
   \item For polynomials of degree 1 on the interval $[a,b]$ it holds that:
\[
f(x) \equiv x \implies \displaystyle \mathcal{E}_{\mathcal{M}_{a}^{b}}\,[f]   = \left[\int_a^b x\,dx\right]-\left[\left(b-a\right)\cdot\left(\frac{a+b}{2}\right)\right] = \frac{b^2-a^2}{2}-\frac{b^2-a^2}{2} = 0.
\]
\end{itemize}
\noindent Hence, for polynomial functions with deg$P\leq1$, the midpoint rule is exact.\newline

\noindent Note that this result aligns with Theorem 4.7 which gives us:
\[
\displaystyle
\left[\left|\int_{a}^{b} f(x)\,dx - \mathcal{M}_{a}^{b}[f]\right|\right]_{f\in P_1}
\leq \frac{1}{24}\,\|f''\|_{L^\infty}\,(b-a)^{3} = \frac{1}{24}\cdot0\cdot (b-a)^3 = 0.
\]

\item To prove that the midpoint rule is not exact on polynomials of degree two, we observe that given:
\[
\displaystyle \mathcal{M}_{a}^{b}[f] := (b-a)\,f\left(\frac{a+b}{2}\right)
\]
  \noindent For polynomials of degree $2$ on the interval $[a,b]$ it holds that:
\[
f(x)\equiv x^2 \implies \displaystyle \mathcal{E}_{\mathcal{M}_{a}^{b}}\,[f]  = \left[\int_a^b x^2\,dx\right]-\left[\left(b-a\right)\left(\frac{a+b}{2}\right)^2\right] = \frac{b^3-a^3}{3} - \frac{(b-a)(a+b)^2}{4}.
\]
Exactness would require:
\[
\frac{b^3-a^3}{3} \equiv  \frac{(b-a)(a+b)^2}{4}.
\]

Which clearly does not hold for all values of $a,b \in \mathbb{R}$. Hence the midpoint rule is not not exact on polynomial functions $P$ with deg$P = 2$.\newline

\noindent Note that this result aligns with Theorem 4.7 which gives us:
\[
\displaystyle
\left[\left|\int_{a}^{b} f(x)\,dx - \mathcal{M}_{a}^{b}[f]\right|\right]_{f\in P_2}
\leq \frac{1}{24}\,\|f''\|_{L^\infty}\,(b-a)^{3} = \frac{\left(b-a\right)^3}{12} \neq 0 \, \, ;\, b\neq a.
\]
\end{enumerate}

\noindent\textbf{Assessed exercise 2.3}
We saw in lectures that, by using a Lagrange interpolation based quadrature with \(n+1\) points, we are able to exactly integrate polynomials up to order \(n\).
However it was seen that with the three points \(a,\ \frac{a+b}{2}, \text{ and } b\), that is \(n=2\), it was possible to integrate a third order polynomial. This is because the points were conveniently chosen.

Here we will work out how to choose four points so that we can integrate fourth order polynomials.
Assuming that the points are given by \(a\), the symmetric points \(\frac{a+b}{2}\pm r\), and \(b\), where \(r\in\left(0,\frac{b-a}{2}\right)\) is to be determined.

\begin{enumerate}
  \item Calculate the Lagrange polynomials associated to this.
  \item Given that the weights associated to the quadrature are given by
  \[
  \omega_0=\omega_3=\frac{(b-a)^3-12r^2(b-a)}{6\left((b-a)^2-4r^2\right)},
  \qquad
  \omega_1=\omega_2=\frac{(b-a)^3}{3\left((b-a)^2-4r^2\right)},
  \tag{1}
  \]
  show that the quadrature is exact on polynomials of order four or less when $  r=\frac{b-a}{2\sqrt{5}}$.
  \item Check if it is exact on fifth order polynomials; for simplicity, set \(a=-1\) and \(b=1\).
\end{enumerate}

\noindent \textbf{SOLUTION:}

\begin{enumerate}
    \item To calculate the Lagrange polynomials, we begin by denoting:
    \[
\displaystyle
x_0=a,\qquad
x_1=\frac{a+b}{2}-r,\qquad
x_2=\frac{a+b}{2}+r,\qquad
x_3=b.
\]
\noindent Then we apply the formula to compute the Lagrange polynomials: 
\[
\displaystyle l_i(x) =\prod_{\substack{j=0 \\ j \neq i}}^{n} \frac{x - x_j}{x_i - x_j}.
\]
Which yields:
\begin{align*}
l_0(x) &= \begin{aligned}[t]
\prod_{j\ne 0}\frac{x-x_j}{x_0-x_j}
&= \frac{x-\left(\frac{a+b}{2}-r\right)}{a-\left(\frac{a+b}{2}-r\right)}
\cdot
\frac{x-\left(\frac{a+b}{2}+r\right)}{a-\left(\frac{a+b}{2}+r\right)}
\cdot
\frac{x-b}{a-b}\\
&= \frac{x-\left(\frac{a+b}{2}-r\right)}{\left(\frac{a-b}{2}+r\right)}
\cdot
\frac{x-\left(\frac{a+b}{2}+r\right)}{\left(\frac{a-b}{2}-r\right)}
\cdot
\frac{x-b}{a-b}\\
&= \frac{
\left(x-\left(\frac{a+b}{2}-r\right)\right)
\left(x-\left(\frac{a+b}{2}+r\right)\right)
(x-b)
}{
(a-b)\left(\left(\frac{a-b}{2}\right)^{2}-r^{2}\right)
}
\end{aligned}
\\[8pt]
l_1(x) &= \begin{aligned}[t]
\prod_{j\ne 1}\frac{x-x_j}{x_1-x_j}
&= \frac{x-a}{\left(\frac{a+b}{2}-r\right)-a}
\cdot
\frac{x-\left(\frac{a+b}{2}+r\right)}{\left(\frac{a+b}{2}-r\right)-\left(\frac{a+b}{2}+r\right)}
\cdot
\frac{x-b}{\left(\frac{a+b}{2}-r\right)-b}\\
&= \frac{x-a}{\left(\frac{b-a}{2}-r\right)}
\cdot
\frac{x-\left(\frac{a+b}{2}+r\right)}{-2r}
\cdot
\frac{x-b}{\left(\frac{a-b}{2}-r\right)}\\
&= \frac{
(x-a)\left(x-\left(\frac{a+b}{2}+r\right)\right)(x-b)
}{
2r\left(\left(\frac{b-a}{2}\right)^{2}-r^{2}\right)
}
\end{aligned}
\\[8pt]
l_2(x) &= \begin{aligned}[t]
\prod_{j\ne 2}\frac{x-x_j}{x_2-x_j}
&= \frac{x-a}{\left(\frac{a+b}{2}+r\right)-a}
\cdot
\frac{x-\left(\frac{a+b}{2}-r\right)}{\left(\frac{a+b}{2}+r\right)-\left(\frac{a+b}{2}-r\right)}
\cdot
\frac{x-b}{\left(\frac{a+b}{2}+r\right)-b}\\
&= \frac{x-a}{\left(\frac{b-a}{2}+r\right)}
\cdot
\frac{x-\left(\frac{a+b}{2}-r\right)}{2r}
\cdot
\frac{x-b}{\left(\frac{a-b}{2}+r\right)}\\
&= -\frac{
(x-a)\left(x-\left(\frac{a+b}{2}-r\right)\right)(x-b)
}{
2r\left(\left(\frac{b-a}{2}\right)^{2}-r^{2}\right)
}
\end{aligned}
\\[8pt]
l_3(x) &= \begin{aligned}[t]
\prod_{j\ne 3}\frac{x-x_j}{x_3-x_j}
&= \frac{x-a}{b-a}
\cdot
\frac{x-\left(\frac{a+b}{2}-r\right)}{b-\left(\frac{a+b}{2}-r\right)}
\cdot
\frac{x-\left(\frac{a+b}{2}+r\right)}{b-\left(\frac{a+b}{2}+r\right)}\\
&= \frac{x-a}{b-a}
\cdot
\frac{x-\left(\frac{a+b}{2}-r\right)}{\left(\frac{b-a}{2}+r\right)}
\cdot
\frac{x-\left(\frac{a+b}{2}+r\right)}{\left(\frac{b-a}{2}-r\right)}\\
&= \frac{
(x-a)\left(x-\left(\frac{a+b}{2}-r\right)\right)\left(x-\left(\frac{a+b}{2}+r\right)\right)
}{
(b-a)\left(\left(\frac{b-a}{2}\right)^{2}-r^{2}\right)
}
\end{aligned}
\end{align*}

\noindent Note as one would expect, integrating our above obtained Lagrange Polynomials: \newline$\displaystyle \int_a^b l_i(x)dx \,;\,i=0,1,2,3$ in WolframAlpha returns the provided values of $\omega_i:$
\[
\omega_0=\omega_3=\frac{(b-a)^3-12r^2(b-a)}{6\left((b-a)^2-4r^2\right)},
\qquad
\omega_1=\omega_2=\frac{(b-a)^3}{3\left((b-a)^2-4r^2\right)},
\]

\item To prove that the quadrature is exact on polynomials of order four or less with $r = \frac{b-a}{2\sqrt{5}}$ we need to show that for $k=0,1,2,3,4$:
\[
\int_a^b x^k\,dx
=\int_a^b \sum_{i=0}^3 x_i^k\,l_i(x)\,dx
=\sum_{i=0}^3 x_i^k \int_a^b l_i(x)\,dx
=\sum_{i=0}^3 \omega_i x_i^k = Q[f].
\]
With
\[
x_0=a,\qquad x_1=\frac{a+b}{2}-r,\qquad x_2=\frac{a+b}{2}+r,\qquad x_3=b
\]
\begin{itemize}
    \item 
    For $k=0$ we have:
\[
\int_a^b 1\,dx=b-a,
\qquad
Q[1]=\omega_0+\omega_1+\omega_2+\omega_3=2\omega_0+2\omega_1.
\]
Evaluating $\omega_0+\omega_1$ yields:
\begin{align*}
\omega_0+\omega_1
&=\frac{(b-a)^3-12r^2(b-a)}{6\left((b-a)^2-4r^2\right)}
  +\frac{(b-a)^3}{3\left((b-a)^2-4r^2\right)}\\
&=\frac{(b-a)^3-12r^2(b-a)+2(b-a)^3}{6\left((b-a)^2-4r^2\right)}\\
&=\frac{3(b-a)^3-12r^2(b-a)}{6\left((b-a)^2-4r^2\right)}
=\frac{(b-a)\left(3(b-a)^2-12r^2\right)}{6\left((b-a)^2-4r^2\right)}\\
&=\frac{(b-a)\cdot 3\left((b-a)^2-4r^2\right)}{6\left((b-a)^2-4r^2\right)}
=\frac{b-a}{2}.
\end{align*}
Hence \(Q[1]=2(\omega_0+\omega_1)=b-a=\int_a^b 1\,dx\).\newline 
\item For $k=1$ we have:
\[
\int_a^b x\,dx=\frac{b^2-a^2}{2}.
\]
\begin{align*}
Q[x] &= \omega_0x_0 + \omega_1x_1 + \omega_2x_2+\omega_3x_3 \\ 
&=\omega_0 a+\omega_1\left(\frac{a+b}{2}-r\right)+\omega_2\left(\frac{a+b}{2}+r\right)+\omega_3 b\\
&=\omega_0 a+\omega_1\left(\frac{a+b}{2}-r\right)+\omega_1\left(\frac{a+b}{2}+r\right)+\omega_0 b\\
&=(\omega_0+\omega_1)(a+b)
=\underbrace{\left(\frac{b-a}{2}\right)}_{\text{As per $k=0$}}(a+b)
=\frac{b^2-a^2}{2}
\end{align*}
Hence \(Q[x]= \omega_0x_0 + \omega_1x_1 + \omega_2x_2+\omega_3x_3=\frac{b^2-a^2}{2}=\int_a^b x\,dx\).\newline

\item For $k=2$ we have:
\[
\int_a^b x^2\,dx=\frac{b^3-a^3}{3}.
\]

\begin{align*}
Q[x^2] 
&=\omega_0 x_0^2 +\omega_1x_1^2+\omega_2x_2^2+\omega_3 x_3^2\\
&=\omega_0 a^2+\omega_1\left(\frac{a+b}{2}-r\right)^2+\omega_1\left(\frac{a+b}{2}+r\right)^2+\omega_0 b^2\\
&=\omega_0(a^2+b^2)+\omega_1\left(\left(\frac{a+b}{2}-r\right)^2+\left(\frac{a+b}{2}+r\right)^2\right)\\
&=\omega_0(a^2+b^2)+\omega_1\left(2\left(\frac{a+b}{2}\right)^2+2r^2\right)\\
&=\omega_0(a^2+b^2)+\omega_1\left(\frac{(a+b)^2}{2}+2r^2\right)\\
&=\omega_0\left(\frac{(a+b)^2+(b-a)^2}{2}\right)+\omega_1\left(\frac{(a+b)^2}{2}+2r^2\right)\\
&=\frac{(b-a)\left((b-a)^2-12r^2\right)}{6\left((b-a)^2-4r^2\right)}\cdot\frac{(a+b)^2+(b-a)^2}{2} +\frac{(b-a)^3}{3\left((b-a)^2-4r^2\right)}\left(\frac{(a+b)^2}{2}+2r^2\right)\\
&=\frac{(b-a)}{6\left((b-a)^2-4r^2\right)}
\left(
\frac{\left((b-a)^2-12r^2\right)\left((a+b)^2+(b-a)^2\right)}{2}
+2(b-a)^2\left(\frac{(a+b)^2}{2}+2r^2\right)
\right)\\
&=\frac{(b-a)}{12\left((b-a)^2-4r^2\right)}
\left(
\left((b-a)^2-12r^2\right)\left((a+b)^2+(b-a)^2\right)
+2(b-a)^2\left((a+b)^2+4r^2\right)
\right)\\
&=\frac{(b-a)}{12\left((b-a)^2-4r^2\right)}
\left(
\begin{aligned}
&(b-a)^2(a+b)^2+(b-a)^4-12r^2(a+b)^2-12r^2(b-a)^2\\
&\quad+2(b-a)^2(a+b)^2+8(b-a)^2 r^2
\end{aligned}
\right)\\
&=\frac{(b-a)}{12\left((b-a)^2-4r^2\right)}
\left(
3(b-a)^2(a+b)^2+(b-a)^4-12r^2(a+b)^2-4(b-a)^2 r^2
\right)\\
&=\frac{(b-a)}{12\left((b-a)^2-4r^2\right)}
\left(
\left((b-a)^2-4r^2\right)\left(3(a+b)^2+(b-a)^2\right)
\right)\\
&=\frac{(b-a)}{12}\left(3(a+b)^2+(b-a)^2\right)
=\frac{(b-a)}{12}\left(3(a^2+2ab+b^2)+a^2-2ab+b^2\right)\\
&=\frac{(b-a)}{12}\left(4a^2+4ab+4b^2\right)
=\frac{(b-a)}{3}\left(a^2+ab+b^2\right)
=\frac{b^3-a^3}{3}.
\end{align*}


Hence \(Q[x^2]=\omega_0 x_0^2 +\omega_1x_1^2+\omega_2x_2^2+\omega_3 x_3^2=\frac{b^3-a^3}{3}=\int_a^b x^2\,dx\).\newline

\item For $k=3$ we have:
\[
\int_a^b x^3\,dx=\frac{b^4-a^4}{4}.
\]

\begin{align*}
Q[x^3]
&=\omega_0 x_0^3+\omega_1 x_1^3+\omega_2 x_2^3+\omega_3 x_3^3\\
&=\omega_0 a^3+\omega_1\left(\frac{a+b}{2}-r\right)^3+\omega_1\left(\frac{a+b}{2}+r\right)^3+\omega_0 b^3\\
&=\omega_0(a^3+b^3)+\omega_1\left(\left(\frac{a+b}{2}-r\right)^3+\left(\frac{a+b}{2}+r\right)^3\right)\\
&=\omega_0(a^3+b^3)+\omega_1\left(
\begin{aligned}
&\left(\frac{a+b}{2}\right)^3-3\left(\frac{a+b}{2}\right)^2 r+3\left(\frac{a+b}{2}\right) r^2-r^3\\
&\qquad\qquad\qquad\,\,
+\left(\frac{a+b}{2}\right)^3+3\left(\frac{a+b}{2}\right)^2 r+3\left(\frac{a+b}{2}\right) r^2+r^3
\end{aligned}
\right)\\
&=\omega_0(a^3+b^3)+\omega_1\left(
2\left(\frac{a+b}{2}\right)^3+\left(-3+3\right)\left(\frac{a+b}{2}\right)^2 r
+6\left(\frac{a+b}{2}\right) r^2+\left(-1+1\right)r^3
\right)\\
&=\omega_0(a^3+b^3)+\omega_1\left(2\left(\frac{a+b}{2}\right)^3+6\left(\frac{a+b}{2}\right)r^2\right)\\
&=\omega_0(a^3+b^3)+\omega_1\left(\frac{(a+b)^3}{4}+3(a+b)r^2\right)\\
&=\frac{(b-a)\left((b-a)^2-12r^2\right)}{6\left((b-a)^2-4r^2\right)}(a^3+b^3)
+\frac{(b-a)^3}{3\left((b-a)^2-4r^2\right)}\left(\frac{(a+b)^3}{4}+3(a+b)r^2\right)\\
& \text{Let } s=a+b,\ d=b-a,\ \implies a=\frac{s-d}{2},\ b=\frac{s+d}{2}\\
& \implies \displaystyle 
\frac{d\,(d^{2}-12r^{2})}{6\,(d^{2}-4r^{2})}(a^{3}+b^{3})
+\frac{d^{3}}{3\,(d^{2}-4r^{2})}\!\left(\frac{s^{3}}{4}+3sr^{2}\right)\\
& a^3 + b^3 =\frac{(s-d)^3+(s+d)^3}{8}=\frac{2s^3+6sd^2}{8} = \frac{s(s^2+3d^2)}{4}\\
&\implies \frac{d\,(d^{2}-12r^{2})}{6\,(d^{2}-4r^{2})}\cdot\frac{s\,(s^{2}+3d^{2})}{4}
+\frac{d^{3}}{3\,(d^{2}-4r^{2})}\!\left(\frac{s^{3}}{4}+3sr^{2}\right)\\
\displaystyle
&=\frac{s}{24\,(d^{2}-4r^{2})}\left[d\,(d^{2}-12r^{2})(s^{2}+3d^{2})
+8d^{3}\!\left(\frac{s^{2}}{4}+3r^{2}\right)\right]\\
\displaystyle
&=\frac{s}{24\,(d^{2}-4r^{2})}\left[d^{3}s^{2}+3d^{5}-12ds^{2}r^{2}-36d^{3}r^{2}
+2d^{3}s^{2}+24d^{3}r^{2}\right]\\
\displaystyle
&=\frac{s}{24\,(d^{2}-4r^{2})}\left[3d^{3}s^{2}+3d^{5}-12ds^{2}r^{2}-12d^{3}r^{2}\right]\\
\displaystyle
&=\frac{s}{24\,(d^{2}-4r^{2})}\cdot 3d\left[d^{2}s^{2}+d^{4}-4s^{2}r^{2}-4d^{2}r^{2}\right]\\
\displaystyle
&=\frac{s}{24\,(d^{2}-4r^{2})}\cdot 3d\,(s^{2}+d^{2})(d^{2}-4r^{2})\\
\displaystyle
&=\frac{sd\,(s^{2}+d^{2})}{8}
=\frac{(a+b)(b-a)\left((a+b)^{2}+(b-a)^{2}\right)}{8}\\
\displaystyle
&=\frac{(b^{2}-a^{2})\left(a^{2}+b^{2}\right)}{4}
=\frac{b^{4}-a^{4}}{4}.
\end{align*}

Hence \(Q[x^3]=\omega_0 x_0^3 +\omega_1x_1^3+\omega_2x_2^3+\omega_3 x_3^3=\frac{b^4-a^4}{4}=\int_a^b x^3\,dx\).\newline

\item For $k=4$ we have: \(\displaystyle \int_a^b x^4\,dx=\frac{b^5-a^5}{5}\).
\begin{align*}
Q[x^4]
&=\omega_0 x_0^4+\omega_1 x_1^4+\omega_2 x_2^4+\omega_3 x_3^4\\
&=\omega_0 a^4+\omega_1\left(\frac{a+b}{2}-r\right)^{\!4}
+\omega_1\left(\frac{a+b}{2}+r\right)^{\!4}
+\omega_0 b^4\\
&=\omega_0(a^4+b^4)+\omega_1\!\left[\left(\frac{a+b}{2}-r\right)^{\!4}+\left(\frac{a+b}{2}+r\right)^{\!4}\right]\\
&=\omega_0(a^4+b^4)+\omega_1\!\left(2\left(\frac{a+b}{2}\right)^{\!4}+12\left(\frac{a+b}{2}\right)^{\!2}r^{2}+2r^{4}\right)\\
&=\omega_0(a^4+b^4)+\omega_1\!\left(\frac{(a+b)^4}{8}+3(a+b)^2r^{2}+2r^{4}\right)\\[2pt]
&=\frac{(b-a)\left((b-a)^2-12r^2\right)}{6\left((b-a)^2-4r^2\right)}\,(a^4+b^4)
\;+\;\frac{(b-a)^3}{3\left((b-a)^2-4r^2\right)}
\left(\frac{(a+b)^4}{8}+3(a+b)^2r^{2}+2r^{4}\right)\\
& \text{Let } s=a+b,\ d=b-a,\ \implies a=\frac{s-d}{2},\ b=\frac{s+d}{2}\\
&a^4+b^4=\frac{(s-d)^{4}+(s+d)^{4}}{16}
=\frac{s^{4}+6s^{2}d^{2}+d^{4}}{8}\\
&\implies \frac{1}{24(d^{2}-4r^{2})}\left[\frac{1}{2}d(d^{2}-12r^{2})(s^{4}+6s^{2}d^{2}+d^{4}) +d^{3}s^{4}+24d^{3}s^{2}r^{2}+16d^{3}r^{4}\right]\\
&\text{At this point note:}\\
&\displaystyle \frac{b^{5}-a^{5}}{5}=\frac{1}{5}\left[\left(\frac{s+d}{2}\right)^{5}-\left(\frac{s-d}{2}\right)^{5}\right]\\
&\text{And: }\displaystyle (s\pm d)^{5} = s^{5} \pm 5s^{4}d+10s^{3}d^{2}\pm 10s^{2}d^{3}+5sd^{4}\pm d^{5}\\
&\implies \displaystyle (s+d)^{5}-(s-d)^{5}
=2\left(5s^{4}d+10s^{2}d^{3}+d^{5}\right)\\
&\implies \frac{b^{5}-a^{5}}{5} = \frac{1}{80}\left(5ds^4+10d^3s^2+d^5\right)\\
&\text{By WolframAlpha our previous expression evaluates to:}\\
&s^2\left(\frac{d^3}{8}+\frac{ds^2}{16}\right) + d^3\left(\frac{d^2}{48}-\frac{r^2}{6}\right)\\
&=\frac{1}{80}\left(5ds^{4}+10d^{3}s^{2}+d^{5}\right)
+\frac{d^{3}}{120}\,(d^{2}-20r^{2})\\
&=\frac{b^{5}-a^{5}}{5}+\frac{(b-a)^{3}}{120}\left((b-a)^{2}-20r^{2}\right)\\
&\text{Hence to obtain exactness on fourth order polynomials, we require }:\\
&(b-a)^{2}=20r^{2} \implies r=\frac{b-a}{2\sqrt{5}}\\
&\text{I.e. given } r=\frac{b-a}{2\sqrt{5}}, \text{ it holds that: }\\
&Q[x^4]=\omega_0 x_0^4 +\omega_1x_1^4+\omega_2x_2^4+\omega_3 x_3^4=\frac{b^5-a^5}{5}=\int_a^b x^4\,dx
\end{align*}
\end{itemize}
Since we have now shown that the quadrature is exact on the monomials $x^k ;k=0,1,2, 3$ and, provided that $r = \frac{b-a}{2\sqrt{5}}$, also for $k=4$. It therefore holds by the linearity of the integral operator that the quadrature is exact on all polynomials of order four or less. I.e.:
\[
\text{Given }r=\frac{b-a}{2\sqrt{5}}, \, a,b\in \mathbb{R}, f\in \mathcal{P}_4\,:\, \mathcal{E}_{\mathcal{Q}_{a}^{b}}\,[f] = 0
\]


\item To check if the quadrature is exact on fifth order polynomials with $a = -1$ and $b = 1$, we can repeat the checks conducted above for $k=0,1,2,3,4$ to the case $k=5$  with $a = -1$ and $b = 1$.\newline


For $k=5$ we have:
\[
\int_{-1}^1 x^5\,dx=\frac{1-1}{6} = 0.
\]

\begin{align*}
Q[x^5]
&=\omega_0 x_0^5+\omega_1 x_1^5+\omega_2 x_2^5+\omega_3 x_3^5\\
&=\omega_0 a^5+\omega_1\left(\frac{a+b}{2}-r\right)^{\!5}
+\omega_1\left(\frac{a+b}{2}+r\right)^{\!5}
+\omega_0 b^5\\
&=\omega_0(1^5+(-1)^5)+\omega_1\!\left[\left(\frac{-1+1}{2}-r\right)^{\!5}+\left(\frac{-1+1}{2}+r\right)^{\!5}\right]\\
&=0 +\omega_1\!\left[\left(0-r\right)^{\!5}+\left(0+r\right)^{\!5}\right]\\
&=\omega_1\!\left[-r^{\!5}+r^{\!5}\right] = 0
\end{align*}
Hence \(Q[x^5]=\omega_0 x_0^5 +\omega_1x_1^5+\omega_2x_2^5+\omega_3 x_3^5=\frac{1^6-(-1)^6}{6}=0=\int_{-1}^1 x^5\,dx\).\newline

\noindent Meaning the quadrature is indeed exact on fifth order polynomials. (Since we can rescale the given interval $[-1,1]$ and achieve the result generally $\forall \, I \in \mathbb{R}$).
\end{enumerate}
\hrulefill
\end{document}


